\section{Introduction}
\label{sec:introduction}

The head of operational risk of a bank owns a number that nobody can check. The economic capital and the internal capital adequacy
assessment of the bank rest on the 99.9 percent quantile of the annual operational loss, and the loss distribution approach estimates it from a
few years of internal losses recorded above a collection threshold \citep{shevchenko2010,cruz2015}. The regulator has already given up
on the same exercise as a Pillar~1 model: the Basel Committee removed the advanced measurement approach because ``the Advanced Measurement Approach's (AMA) inherent
complexity, and the lack of comparability arising from a wide range of internal modelling practices, have exacerbated variability in
risk-weighted asset calculations'' \citep{bcbsd355}, and the final standard replaced the approaches by one standardised formula
\citep{bcbsd424,ope25}. The loss distribution approach nevertheless remains the model of economic capital and the object that a
validator is asked to judge \citep{bcbs196}.

The shape of the data explains why the judgement is hard. In the 2008 Loss Data Collection Exercise, 119 of the 121 participating banks submitted 10.6 million
internal losses worth EUR 59.6 billion; the 171,882 losses of EUR 20,000 or more were 1.6 percent of the count and 90 percent of the amount,
and the twenty largest losses alone were 29.5 percent of the total \citep{bcbs160}. A quantile at the level 0.999 read from five years of such
data is a quantile of a few dozen large losses, extrapolated by a tail model far beyond the largest observation. Its estimator inherits the instability of the tail index that governs the extrapolation. A validator can recompute the number and can test the fit of the
severity, but cannot say how wrong the number is, because the true quantile of a bank is never observed.

\paragraph{The gap.} The published literature on the loss distribution approach supplies the estimators, the aggregation algorithms and the warnings.
Aggregation by recursion and by Fourier transform is exact on a discretised severity \citep{panjer1981,shevchenko2010}; the truncation at the collection threshold biases the capital if it is
ignored \citep{mignola2007,cope2009}; tails of infinite mean break the usual risk calculus \citep{neslehova2006,chavez2006}; maximum likelihood is fragile and robust estimators are recommended
\citep{opdyke2012}; and the standardised approach has been criticised for insensitivity and instability \citep{peters2016,mignola2016}. What these sources do not report is a measurement:
the dispersion of the capital of each estimator over independent histories, against a quantile that is known, on the same data for all
estimators. The reasons given by the Basel Committee for dropping the advanced approach are complexity and lack of comparability
\citep{bcbsd355}, not a measured error. The academic critique argues the instability of the standardised approach by analysis and by examples
\citep{peters2016}. In the sources surveyed in
\cref{sec:related-work} we found no study that measures how the 99.9 percent capital of the loss distribution approach disperses against a known truth over independent
histories; that absence is a statement about the sources we searched, not a proof of absence.

\paragraph{What we do.} We propose a measurement protocol and apply it. The protocol draws losses from a generator whose annual loss distribution is known, computes the true 99.9 percent
quantile once by three independent aggregations, and then draws 1{,}000 independent histories on disjoint seeds in each cell of a design with four tail indices (0.6, 0.8, 0.96 and 1.1),
two history lengths (5 and 10 years) and two collection thresholds (EUR 10,000 and 20,000). Each history goes through a ladder of ten estimation rungs and through the standardised approach,
all computed by one engine (\texttt{riskvalidation} 1.1.0, \cref{sec:pm-implementation}), and each capital is divided by the true quantile. The protocol reports, per rung and per cell, the median of the ratio, its interquartile range,
the share of histories that understate the truth, the number of histories on which the rung returns a capital at all, and the share of histories whose fitted tail index reaches one,
that is, that would conclude that the mean does not exist. Each statistic carries a bootstrap or Wilson interval over the histories \citep{efron1979,wilson1927}.

\paragraph{Contributions.}
\begin{enumerate}
\item A measurement protocol for the instability of the operational-risk capital on independent histories of known truth, with its metrics, its seeds and a measured precision for the number of histories it draws (\cref{sec:proposed-method,sec:res-truth}).
\item A measurement on \nHist{} histories: the interquartile range of the splice capital is \spIqrMin{} to \spIqrMax{} times the true quantile, the capital stability test is red in every cell where a rung returns a capital (\redTotal{} of \redTotal{} rows, of which \redSix{} of \rowsSix{} on the six-rung ladder),
and the share of histories that conclude an infinite mean at tail index 0.96 is \infLo{} to \infHi{} percent (\cref{sec:results}).
\item A delta-method theory of the single loss approximation fitted by maximum likelihood to a generalised Pareto tail, which predicts the measured dispersion within \thSdMaxErr{} percent in the 16 cells, predicts the share of false infinite-mean conclusions within \thShareMaxErr{} points, and gives the history length that would bring the dispersion under a stated limit (\cref{sec:theory}).
\item An ablation that measures what each estimation component contributes on the same histories, including components that do not help, such as the minimum Cram\'er-von Mises fit of the tail (\cref{sec:ablation}).
\end{enumerate}

\paragraph{Main result.} More data does not make the number stable at the horizons a bank has. Doubling the history from 5 to 10 years reduces the dispersion of the single loss approximation by the factor $\sqrt2$ that the theory gives and no more, and the history
needed to reach an interquartile range of 0.25 times the truth is \yrsMin{} to \yrsMax{} years across the cells. The standardised approach is stable across histories relative to its own level, because it depends on the losses only through their average annual loss, and it does not track the truth: in the cell of 10 years and a threshold of EUR 20,000 its interquartile range is \smaIqrA{}, \smaIqrB{}, \smaIqrC{} and \smaIqrD{} times the true quantile and its median capital is \smaA{}, \smaB{}, \smaC{} and \smaD{} times the true quantile at the four tail indices, in that order (\cref{sec:results}).
The hypothesis stated in the case parameters before the robust fits were run, that the Cram\'er-von Mises and the bias-robust fits beat the splice in spread and in bias, is inconclusive over the 16 cells and refuted at the tail indices 0.96 and 1.1 (\cref{sec:res-hypotheses}).

\paragraph{Scope.} The bank is synthetic and its severity belongs to the generalised Pareto family that the tail estimators assume, which favours them; the lights of the capital stability test are internal limits of the product and not regulatory thresholds. The numbers say nothing about the capital of any real bank,
and \cref{sec:limitations} states what the design cannot show.

\paragraph{Structure.} \Cref{sec:related-work} places the protocol among the published methods, \cref{sec:problem} states the problem and the notation, \cref{sec:proposed-method} describes the protocol and its algorithm, \cref{sec:theory} derives the dispersion and the misclassification share,
\cref{sec:experimental-design} states the hypotheses and the design, \cref{sec:results} reports the main comparison, \cref{sec:ablation} the ablation, and \cref{sec:discussion,sec:limitations,sec:reproducibility,sec:conclusions} close the work.
